\documentclass[10pt,a4paper]{amsart}
\usepackage[margin=19mm]{geometry}
\usepackage{amsmath,amssymb,mathtools}
\usepackage[scr=rsfs,cal=euler]{mathalfa}
\usepackage{xfrac}
\usepackage[unicode,hidelinks,bookmarks=false]{hyperref}
\newcommand{\ie}{i.e.\ }
\newcommand{\cf}{cf.\ }
\newcommand{\resp}{respectively\ }
\newcommand{\Subsection}{\subsection}
\providecommand{\nobreakdash}{\nobreak\hbox{-}}
\setlength{\emergencystretch}{2em}
\multlinegap0pt
\usepackage{bm}
\usepackage{bbm}
\usepackage{dsfont}
\usepackage{xspace}
\usepackage{cancel}
\usepackage{mathtools}
\usepackage{amsthm}
\makeatletter
\@ifundefined{thm}{\newtheorem{thm}{Thm}}{}
\@ifundefined{lem}{\newtheorem{lem}[thm]{Lemma}}{}
\@ifundefined{prop}{\newtheorem{prop}[thm]{Proposition}}{}
\makeatother
\newcommand{\Spec}{\operatorname{Spec}}
\newcommand{\nil}{\operatorname{Nil}}
\newcommand{\rad}{\operatorname{rad}}
\title[The prime spectrum of the talented monoid of a higher-rank g]{The prime spectrum of the talented monoid of a higher-rank graph and applications}
\author{Roozbeh Hazrat, Promit Mukherjee}
\date{Source: arXiv:2609.16632v1 (CC BY 4.0)}
\begin{document}
\maketitle

\noindent\emph{Source and licence.} The prime spectrum of the talented monoid of a higher-rank graph and applications. Version arXiv:2609.16632v1 (CC BY 4.0), distributed under the Creative Commons Attribution 4.0 International licence, as recorded in the arXiv licence metadata for this exact version and shown on the abstract page https://arxiv.org/abs/2609.16632v1. Full rights notice: licenses/arxiv-2609.16632-CC-BY-4.0.txt.\par\smallskip

\noindent\textit{Editorial scope.} Counted unit: the Prop whose source label is \texttt{pro basic facts about basic open sets} in the article (source0.tex lines 563--595, source label \texttt{pro basic facts about basic open sets}), delivered together with the article's own complete original proof of it (source0.tex lines 596--626, one \texttt{proof} environment of 31 lines). No mathematics is written, repaired or re-derived: statement and proof are the article's own text line for line. Required context retained uncounted: lem maximal order ideal is prime (lines 530--532); lem proper ideal contained in a prime ideal (lines 541--543). Every definition, display and label of those lines is retained unchanged. Omitted from this package and not redistributed: 1--529, 533--540, 544--562, 627--1127. Nothing inside a retained excerpt is omitted or altered: every retained body is delivered exactly as the article's lines are written, so no omission receipt of any kind accompanies this package. Citations: the retained excerpts carry no citation command at all, so no bibliography entry of the article is transcribed and no reference is redistributed. Reference adaptation: none was needed and none was made; every reference inside the delivered unit resolves inside it, so no unresolved reference or citation is produced. Printed slips kept literal: no printed word, symbol, hypothesis or number of the counted statement or of its proof was altered. Layout and class: the article's own class is \texttt{amsart} with packages including fontenc, inputenc, amsmath, amssymb, amsfonts, amsthm, amscd, mathrsfs, calligra, bm, bbm, booktabs, tabularx, tikz; the wrapper is the lane's pinned \texttt{amsart} editorial wrapper at 10pt on A4, which restates the theorem-like environments the retained text needs and adds the article's own macro definitions that the retained text needs (\texttt{\textbackslash Spec}, \texttt{\textbackslash nil}, \texttt{\textbackslash rad}), with extra wrapper packages (bm, bbm, dsfont, xspace, cancel, mathtools). The delivered numbering is the unit's own. Assets: no retained span contains a figure environment, a graphics-inclusion command, a PDF-page inclusion, a table or a TikZ picture; no image file is copied into this package, the package declares no asset dependency and no image file is redistributed with it. Licence: version arXiv:2609.16632v1 of this article is distributed under the Creative Commons Attribution 4.0 International licence, as recorded in the arXiv licence metadata for this exact version and shown on the abstract page https://arxiv.org/abs/2609.16632v1. A full rights notice is delivered at licenses/arxiv-2609.16632-CC-BY-4.0.txt. Characters: every retained excerpt is pure ASCII, so the delivered engine, XeLaTeX with its default Latin Modern text font, reports no missing character.
\begin{lem}\label{lem maximal order ideal is prime}
Let $\Gamma$ be a group and $M$ a $\Gamma$-monoid with the refinement property. Then every maximal $\Gamma$-order ideal of $M$ is prime.    
\end{lem}

\begin{lem}\label{lem proper ideal contained in a prime ideal}
Let $\Gamma$ be a group and $M$ a $\Gamma$-monoid with the refinement property. Then for every proper $\Gamma$-order ideal $J$ of $M$ and $x\notin J$, there is a prime $\Gamma$-order ideal $P$ such that $J\subseteq P$ and $x\notin P$. In particular, if $M$ is conical and $M\neq \{0\}$, then $\Spec_\Gamma(M)\neq \emptyset$.    
\end{lem}

\begin{prop}\label{pro basic facts about basic open sets}
Let $\Gamma$ be a group and $M$ a conical $\Gamma$-monoid such that $M\neq \{0\}$. Then the following hold.

$(i)$ $D(x)=\emptyset$ if and only if $x\in \nil(M)$.

$(ii)$ $D(x)\cap D(y)=D(\langle x\rangle \cap \langle y\rangle)$ for $x,y\in M$.

$(iii)$ $\overline{\{P\}}=V(P)$ for any $P\in \Spec_\Gamma(M)$.

$(iv)$ $\Spec_\Gamma(M)$ is $T_0$.

$(v)$ $V(J)=V(\rad(J))$ for any $\Gamma$-order ideal $J$ of $M$.

$(vi)$ For any $\Gamma$-order ideal $J$, $\rad(J)$ is prime if and only if $V(J)$ is a nonempty irreducible subset.

$(vii)$ $\Spec_\Gamma(M)$ is sober, i.e., every nonempty irreducible set has a unique generic point.

$(viii)$ $\Spec_\Gamma(M)$ is spectral in the sense of Hofmann--Keimel.

Furthermore, if $M$ is a refinement monoid, then 

$(1)$ $D(x)=\Spec_\Gamma(M)$ if and only if $x$ is an order unit in $M$;

$(2)$ for any $P\in \Spec_\Gamma(M)$, $\{P\}$ is closed if and only if $P$ is maximal;

$(3)$ $J=\rad(J)$ for any $\Gamma$-order ideal of $M$, in particular, $\nil(M)=\{0\}$;

$(4)$ $D(x)=\emptyset$ if and only if $x=0$;

$(5)$ $D(x)$ is compact for all $x\in M$, in particular, $\Spec_\Gamma(M)$ is compact. 

If, in addition, each $\Gamma$-order ideal of $M$ is finitely generated, then $\Spec_\Gamma(M)$ is spectral in the sense of Hochster. 
\end{prop}

\begin{proof}
$(i)$ This is obvious from the definition of nilradical of $M$. 

$(ii)$ This follows from the definition of a prime $\Gamma$-order ideal.

$(iii)$ Obviously, $V(P)$ is a closed set containing $P$. Now if $V(J)$ is any closed set containing $P$, then $J\subseteq P$, whence $V(P)\subseteq V(J)$. Thus, $V(P)=\overline{\{P\}}$.

$(iv)$ Let $P,Q\in \Spec_\Gamma(M)$ be such that $P\neq Q$. Then $V(P)\neq V(Q)$ which implies $\overline{\{P\}}\neq \overline{\{Q\}}$ by $(iii)$. So $\Spec_\Gamma(M)$ is $T_0$. 

$(v)$ If $P\in V(J)$, then $\rad(J)=\displaystyle{\bigcap_{Q\in V(J)}}Q\subseteq P$. Hence $P\in V(\rad(J))$. Since $J\subseteq \rad(J)$, the other inclusion is obvious. 

$(vi)$ Assume that $V(J)$ is a nonempty irreducible subset. Suppose $J_1$, $J_2$ are $\Gamma$-order ideals such that $J_1\cap J_2\subseteq \rad(J)$. Let $P\in V(J)$. Then $J_1\cap J_2\subseteq \rad(J)\subseteq P$. By primeness of $P$, $J_1\subseteq P$ or $J_2\subseteq P$, whence $P\in V(J_1)\cup V(J_2)$. Therefore, $V(J)\subseteq V(J_1)\cup V(J_2)$. Since $V(J)$ is irreducible, $V(J)\subseteq V(J_1)$ or $V(J)\subseteq V(J_2)$. If $V(J)\subseteq V(J_1)$, then $J_1\subseteq P$ for all $P\in V(J)$, whence $J_1\subseteq \rad(J)$. Similarly, if $V(J)\subseteq V(J_2)$, then $J_2\subseteq \rad(J)$. Hence, $\rad(J)$ is a prime $\Gamma$-order ideal. Now suppose $\rad(J)$ is prime. Then $\rad(J)\neq M$ and so $V(J)\neq \emptyset$. To show that $V(J)$ is irreducible, let $I_1,I_2$ be $\Gamma$-order ideals such that $V(J)=V(I_1)\cup V(I_2)$. Then $V(J)=V(I_1\cap I_2)$. Since $\rad(J)$ is a prime $\Gamma$-order ideal containing $J$, $I_1\cap I_2\subseteq \rad(J)$. The primeness of $\rad(J)$ now implies that $I_1\subseteq \rad(J)$ or $I_2\subseteq \rad(J)$. If the first one happens, then for any $P\in V(J)$, $I_1\subseteq P$ and so $P\in V(I_1)$. Hence, $V(J)=V(I_1)$. Similarly, if $I_2\subseteq \rad(J)$, then $V(J)=V(I_2)$. This shows that $V(J)$ is irreducible.  

$(vii)$ Let $X$ be any nonempty irreducible subset of $\Spec_\Gamma(M)$. Then $X=V(J)$ for some $\Gamma$-order ideal $J$ of $M$ such that $\rad(J)$ is prime. Now using $(iii)$ and $(v)$, we have $V(J)=V(\rad(J))=\overline{\{J\}}$, showing that $\rad(J)$ is a generic point for $V(J)$. The uniqueness of the generic points is clear since for any two prime $\Gamma$-order ideals $P, Q$, $V (P ) = V (Q)$ implies $P = Q$.

$(viii)$ This follows from $(iv)$ and $(vii)$. 

Now we assume that $M$ has the refinement property.

$(1)$ The `if' direction is evident as no prime $\Gamma$-order ideal can contain an order unit. Now assume that $D(x)=\Spec_\Gamma(M)$. This implies $\langle x\rangle \nsubseteq P$ for every prime $\Gamma$-order ideal $P$. Then $\langle x\rangle=M$, otherwise if it is a proper $\Gamma$-order ideal, then it would be contained in a prime $\Gamma$-order ideal of $M$ by Lemma \ref{lem proper ideal contained in a prime ideal}. It follows that for any $a\in M$, there is some $\alpha\in \mathbb{N}[\Gamma]$ such that $a\le~^{\alpha}x$. Hence, $x$ is an order unit of $M$.

$(2)$ If $P$ is a maximal $\Gamma$-order ideal, then it is prime by Lemma \ref{lem maximal order ideal is prime} and it is clear that $V(P)=\{P\}$. Therefore, $\{P\}$ is closed. Conversely, if $\{P\}$ is closed, then $\{P\}=V(P)$ by $(iii)$. Let $P\nsubseteq I\subseteq M$ be a $\Gamma$-order ideal. If $I\neq M$, then by Lemma \ref{lem proper ideal contained in a prime ideal}, there is a prime $\Gamma$-order ideal $K$ of $M$ such that $I\subseteq K$. But then $K\in V(P)$ and $K\neq P$ which is a contradiction. Therefore, $I=M$ proving that $P$ is maximal. 

$(3)$ Let $J$ be any $\Gamma$-order ideal of $M$. The inclusion $J\subseteq \rad(J)$ is obvious. Since $\rad(M)=M$, we may assume that $J$ is proper. But then the result follows from Lemma \ref{lem proper ideal contained in a prime ideal}. Since $\rad(\{0\})=\nil(M)$, it follows that $\nil(M)=\{0\}$. 

$(4)$ Follows from $(i)$ and $(3)$.

$(5)$ Let $x\in M$. Since the sets $D(y)$, $y\in M$ form a basis for the Zariski topology on $\Spec_\Gamma(M)$, it suffices to check that any open cover of $D(x)$ consisting of the basic open sets of the form $D(y)$ has a finite subcover. So let $\Delta$ be an index set and $\mathcal{U}=\{D(x_\nu)~|~\nu\in \Delta,x_\nu\in M\}$ an open cover of $D(x)$. Let $J=\langle \{x_\nu~|~\nu\in \Delta\}\rangle$. Let $P\in V(J)$. If $x\notin P$, then $P\in D(x)$ and consequently, $P\in D(x_\nu)$ for some $\nu\in \Delta$. But then $J\nsubseteq P$ which is not the case. Thus $x\in P$ for all $P\in V(J)$ and so $x\in \rad(J)=J$ by $(3)$. It follows that $x\le \displaystyle{\sum_{i=1}^{t}}~^{\gamma_i}x_{\nu_i}$ for some $t\in \mathbb{N}$ and $\gamma_i\in \mathbb{N}[\Gamma]$; $i=1,2,\ldots,t$. Now for any $Q\in D(x)$, we must have $j\in \{1,2,\ldots,t\}$ such that $x_{\nu_j}\notin Q$, i.e., $Q\in D(x_{\nu_j})$. Therefore, $D(x)\subseteq \displaystyle{\bigcup_{i=1}^{t}}~D(x_{\nu_i})$ and we have extracted a finite subcover for $D(x)$ from $\mathcal{G}$. Since $\Spec_\Gamma(M)=D(\epsilon)$, $\Spec_\Gamma(M)$ is also compact.

For the final statement, assume that every $\Gamma$-order ideal of $M$ is finitely generated. Note that $\Spec_\Gamma(M)$ is already spectral in the sense of Hofmann--Keimel. It is compact by $(5)$. Also, by the same part, each member of the base $\mathcal{B}=\{D(x)~|~x\in M\}$ is compact. Let $x,y\in M$. Then by $(ii)$, $D(x)\cap D(y)=D\left(\langle x\rangle \cap \langle y\rangle\right)$. By the assumption, there exist $n\in \mathbb{N}$ and $z_i\in \langle x\rangle \cap \langle y\rangle$, $i=1,2,\ldots,n$ such that $\langle x\rangle \cap \langle y\rangle=\langle \{z_i~|~i=1,2,\ldots,n\}\rangle=\langle \displaystyle{\sum_{i=1}^{n}}~z_i\rangle$. It is then easy to see that \[D(x)\cap D(y)=D\left(\langle x\rangle \cap \langle y\rangle\right)=D\left(\displaystyle{\sum_{i=1}^{n}}~z_i\right).\] Therefore, the base $\mathcal{B}$ is closed under finite intersection. Hence, $\Spec_\Gamma(M)$ is spectral in the sense of Hochster. 
\end{proof}

\end{document}
